\section*{ЗАКЛЮЧЕНИЕ}
\phantomsection
\addcontentsline{toc}{section}{Заключение}

В ходе выполнения курсовой работы проведено численное исследование влияния дискретного микрорельефа на рабочие характеристики радиального подшипника скольжения. Решены следующие задачи:

\begin{enumerate}
    \item Выполнен аналитический обзор основ гидродинамической теории смазки и опубликованных результатов по лазерному текстурированию трибологических поверхностей. Показано, что положительное влияние микрорельефа обусловлено преимущественно микрогидродинамическим эффектом, возникающим при обтекании локальных углублений смазочным потоком.

    \item Реализована методика численного решения уравнения Рейнольдса для подшипника конечной длины на основе метода конечных разностей с итерационной схемой Гаусса--Зейделя и верхней релаксацией. Вычисления проведены на расчётной сетке $500 \times 500$ узлов при критерии сходимости $\delta < 10^{-5}$.

    \item Определены поля давления и интегральные параметры (нагрузочная способность $F$, коэффициент трения $\mu$, расход смазки $Q$) для гладкого подшипника и подшипника со сферическими углублениями (тип~5) при эксцентриситетах $\varepsilon = 0{,}05 - 0{,}80$. Геометрия подшипника (тип~B): $R = 50$~мм, $c = 70$~мкм, $L = 80$~мм.

    \item Проведено сопоставление гладкого и текстурированного подшипников. При рабочем эксцентриситете $\varepsilon = 0{,}6$ установлены следующие изменения:
    \begin{itemize}
        \item нагрузочная способность возросла на $4{,}85$\%: с $F = 10685{,}6$~Н до $F = 11203{,}4$~Н;
        \item коэффициент трения уменьшился на $5{,}34$\%: с $\mu = 0{,}01507$ до $\mu = 0{,}01426$;
        \item расход смазки увеличился на $0{,}45$\%: с $Q = 0{,}6739$~л/с до $Q = 0{,}6770$~л/с.
    \end{itemize}
\end{enumerate}

Результаты моделирования подтверждают, что нанесение сферических углублений с параметрами $a = b = 2{,}3$~мм, $h_p = 12$~мкм ($h_p/c = 0{,}2$) приводит к улучшению трибологических характеристик подшипника. Одновременный рост нагрузочной способности и снижение коэффициента трения при незначительном увеличении расхода смазки свидетельствуют о практической целесообразности данного типа текстурирования.

Профиль сферической шапки обладает симметричной кривизной, что обеспечивает равномерное распределение дополнительного давления вокруг центра углубления без резких перепадов потока. Данная особенность делает сферический микрорельеф пригодным для подшипников, эксплуатируемых в стационарных режимах при умеренных и высоких эксцентриситетах.

На основании выполненного анализа рекомендуется использование сферического микрорельефа типа~5 для подшипников скольжения нефтегазового оборудования, функционирующих в условиях жидкостного трения. Наибольший эффект ожидается в узлах с повышенной нагрузкой ($\varepsilon > 0{,}5$), где прирост несущей способности и уменьшение трения выражены особенно отчётливо.

